\section{The cubic relative relation count at the transported local component}
\label{sec:cubic-box-presentation}

Retain AC-G49, the old source \texttt{CohInp}, and the full
quantified \texttt{LemmaP} in
\path{CommutativeAlgebra/FramedTwoCoboundaries.lean}.
The latter hypothesis is expressly retained; the analytic
theorem A alone does not imply it.
All source paths in this section have the same pinned root
\path{Development/AutomorphicArithmetic/AutomorphicVarieties/}.
Let $B_\chi$ denote the universal cubic ring before imposing
normalization, but after imposing the actual killing and
selected local conditions.  Let
$B_\ell=U_\chi.R/P_\chi$ and choose the local weighted
presentation $W_\ell$ constructed explicitly in the first
step below.  The transported chart of AC-G49-inert
is part of the hypotheses here.

\begin{theorem}[AC-G50: the exact cubic \texttt{BoxPres} relation count]
\label{AC-G50}
There is a finite relative homogeneous presentation of
$B_\chi$ over the actual $B_\ell$ of the following form:
\[
 B_\chi\simeq
 O[[X_s\ (s\in\Sigma_\ell),Y_t\ (t\in T)]]
 \big/\bigl(I_\ell,(f_u)_{u\in U}\bigr).                    \tag{G50.1}
\]
Here $I_\ell$ is the kernel of the chosen local presentation,
all coordinate weights are in $[-2,2]$, the equations $f_u$
are homogeneous, and, with
$a_i=\#\{t:w(Y_t)=i\}$,
\[
 \#\{u:w(f_u)=i\}=a_i+b_i,
 \qquad b_i=\begin{cases}1,&|i|\le1,\\0,&|i|>1.
                       \end{cases}                         \tag{G50.2}
\]
The new coordinates are actual embedded matrix-entry
differences from $r_\chi$, and all universal matrix entries
are homogeneous with their usual weights.
This is every field of \texttt{BoxPres}, with the explicitly
constructed cubic seed and local object substituted.
It is not an inference from dimension counting alone.
\end{theorem}
\begin{proof}
\proofstep{The actual local weighted presentation.}
Let $u$ be the representation of the local Galois group
over $B_\ell$ induced by the universal local representation.
For each local element $g$ and pair of embedded indices
$(a,b)$, use the coordinate
\[
 e_{g,a,b}=\operatorname{emb}(u(g))_{ab}
       -\operatorname{emb}(r_\chi(g))_{ab}.
\]
It lies in the maximal ideal, because the residual matrices
agree.  These coordinates determine every homomorphism out
of $B_\ell$: equality of all of them gives equality of the
local framed lifts, hence equality of their maps out of
$U_\chi.R$ by local universality, and then equality on its
surjective quotient $B_\ell$.

For every source test ring and every unit $t\equiv1$, strict
conjugation by $T=\operatorname{diag}(t,1,t^{-1})$ preserves
the actual local condition $P_\chi$.  Indeed $T$ commutes
with the fixed $C$, and after conjugating by $C^{-1}$ this
is precisely the old strict conjugation preserving $I_t$.
The embedded basepoint has no nonzero entry between
unequal weights, so this operation scales $e_{g,a,b}$ by
$t^{w_6(a)-w_6(b)}$.  The general
\texttt{coord\_grading} theorem therefore applies directly
to $B_\ell$.  It chooses finitely many of these actual
coordinates, constructs their surjective homogeneous
presentation $W_\ell$, and puts every $e_{g,a,b}$ in its
stated weight part.  Adding the basepoint entry preserves
that weight: a nonzero basepoint entry has weight zero.
Thus every embedded entry of $u$ is homogeneous, all
chosen coordinates are maximal-ideal elements, and they
are actual entry differences.  These are exactly both
fields of \texttt{IsLocPres} in
\path{UnitaryThreeGroups/IsLocPres.lean}.
This argument constructs the required local presentation
separately; it does not reuse the global presentation of
$R_\chi$ as though it were a presentation of $B_\ell$.

\proofstep{The unchanged six-dimensional form problem.}
Use the exact source symmetric form and block data
\[
                \mathcal J=\begin{pmatrix}0&J\\J&0\end{pmatrix},
          \qquad w_6=(1,0,-1,1,0,-1).                       \tag{G50.3}
\]
For every $x\in GG_3(A)$ the elementary block calculation
$\operatorname{emb}(x)^t\mathcal J\operatorname{emb}(x)=\mathcal J$
holds.  Its two three-dimensional blocks are preserved
or interchanged according to the multiplier.
The embedded cubic basepoint has nonzero entries only
between equal weights, as in (G49.12).
Thus both $\operatorname{emb}r_\chi$ and its reduction
satisfy every field of \texttt{FormData.Normalises}.
These are precisely the three checks in
\path{UnitaryThreeGroups/BoxDefs.lean},
\texttt{emb\_form}, \texttt{emb\_blk} and
\texttt{blockPairingForm\_normalises\_baseLift}; the first
two statements already quantify over arbitrary $GG_3$
elements, and the third one's weight assertion follows
here from diagonal blocks rather than the special powers
of a single character.

The form-valued lift condition is exactly the $GG_3$ lift
condition for this new seed.  To verify the converse
explicitly, if a lift $R(g)$ differs from the embedded
basepoint by a strict form-preserving block-diagonal matrix,
that matrix has the form
$\operatorname{diag}(A,JA^{-t}J)$.
Multiplying it by the embedded cubic basepoint recovers a
unique $GG_3$ element of multiplier $\eta(g)$.
The faithful embedding proves its group law and its correct
reduction.  This is the matrix proof in
\path{UnitaryThreeGroups/GLiftToLiftInHAndSurjective.lean};
it uses the basepoint only as an actual $GG_3$
homomorphism with the specified multiplier.

\proofstep{The exact away-from-$\ell$ conditions.}
Choose $S$ as in AC-G49 and, for each $v\in S$ away from
$\ell$, let $D_v$ and $I_v$ be a decomposition and inertia
group and put
\[
                       N_v=I_v\cap\ker\overline r_\chi.      \tag{G50.4}
\]
These are closed normal subgroups of $D_v$.
Their indices in inertia are prime to $\ell$.
Indeed the finite residual image of $r_\chi$ has order
dividing $2\cdot3\cdot\operatorname{ord}(\overline\gamma)$,
and $\operatorname{ord}(\overline\gamma)$ divides
$\#k_O^\times$.  The coefficient primes $2,3$ are excluded.
The proof of AC-G49 shows that the characteristic-zero
basepoint kills each $N_v$.

For any one of its finite residual adjoint weight modules,
$N_v$ acts trivially.  The quotient $D_v/I_v$ is procyclic
and has unramified quotients of order $\ell^n$ for every
$n$.  These facts verify all hypotheses of the general
theorem
\path{CommutativeAlgebra/ExistsLocallyConstantDOneEqCocycleOfTameInertia.lean},
\texttt{exists\_locallyConstant\_dOne\_eq\_cocycle\_of\_tame\_inertia}.
It makes every locally constant two-cocycle on $D_v/N_v$
a coboundary.  Applying the equally general
\path{CommutativeAlgebra/KillLoc.lean}, \texttt{kill\_loc},
therefore proves the full \texttt{IsLocCond} for the actual
condition of killing $N_v$, including its infinitesimal
lifting field.  It is not merely a statement about its
characteristic-zero points.

The simultaneous condition of killing every such $N_v$
and all inertia outside $S$ equals killing the closed
normal closure of $\mathcal K_\chi$.  To see both directions,
first note that each $N_v$ is a subset of that union of
inertia kernels.  Conversely, every inertia kernel at a
prime in $S\setminus\{\ell\}$ is a conjugate of one in
(G50.4), and outside $S$ the whole inertia group is killed
in $G_{\mathbb Q,S}$.  Normality and closedness give the
claim.  Thus the frame used below represents the actual
killing condition of AC-G49.

\proofstep{A cyclotomic witness for every cubic weight.}
The literal theorem
\path{UnitaryThreeGroups/ExistsGaloisCyclotomicNegOneTrivialOnEtaPsi.lean},
\texttt{exists\_galois\_cyclotomic\_neg\_one\_trivial\_on\_eta\_psi},
provides an element $g$ with
\[
       \eta(g)=1,\qquad \overline\gamma(g)=1,
                          \qquad\overline\kappa_\ell(g)=-1. \tag{G50.5}
\]
Its use retains the exact old \texttt{GD.eta} hypothesis.
Replace $g$ by $g^3$.  Since $g\in H$ and $q^3=1$,
\[
             \overline r_\chi(g^3)=1,
                      \qquad\overline\kappa_\ell(g^3)=-1.   \tag{G50.6}
\]
This single element acts trivially on every adjoint weight
module.  In particular the cyclotomic witness required in
the generic finite-coefficient \texttt{box\_count} theorem
survives the cubic change; no new class-field existence
hypothesis is inserted.

\proofstep{Compute the actual correction at infinity.}
For each integer $i$, let $M_i$ be the degree-$i$ part of
the residual form Lie algebra, with action by
$\operatorname{emb}\overline r_\chi$.
It can be identified with the corresponding part of
$\mathfrak{gl}_3(k_O)$ using
\[
                       X\longmapsto\operatorname{diag}(X,-JX^tJ).
                                                                  \tag{G50.7}
\]
Complex conjugation has matrix component $I$ in (G49.1),
so its action on $X$ is $X\mapsto-JX^tJ$.
On the degree $2$ and $-2$ lines this is $-1$.
On each of the two-dimensional degree $1$ and $-1$ pieces
it interchanges the two matrix units with a minus sign,
leaving one fixed line.  On the diagonal degree-zero
space its fixed vectors are
$\operatorname{diag}(a,0,-a)$.
Since $\ell\ne2$, this proves
\[
                             \dim_{k_O} M_i^{c=1}=b_i.       \tag{G50.8}
\]
This directly verifies the cubic instance of the calculation
in \path{UnitaryThreeGroups/FinrankKerAdRepComplexConjSubOneEqBoxDef.lean}.

\proofstep{Use the literal generic cohomological count.}
The theorem \texttt{box\_count} in
\path{UnitaryThreeGroups/BoxCount.lean} already quantifies
over an arbitrary finite residual coefficient $M$ and
its representation $\rho$; it is not restricted to a
rational character.
Apply it to each $M_i$, the frame (G50.4), the original
local embedding at $\ell$, and the witness (G50.6).
The preceding steps check its open-kernel, local
prime-to-$\ell$ index and global witness hypotheses.
Its conclusion is
\[
 \dim_{k_O}\operatorname{Ob}(M_i)
       =\dim_{k_O} H^1_c(M_i)+b_i,                          \tag{G50.9}
\]
and both dimensions are finite.  The actual local and
global Euler characteristics and Poitou--Tate input in
\texttt{CohInp} are retained.  In particular no vanishing
of the mixed characteristic-zero Bloch--Kato group has
been used to obtain (G50.9).

\proofstep{Construct the presentation rather than count a hypothetical ring.}
Apply the full \texttt{LemmaP} to (G50.3), the actual
embedded residual and base representations, the frame
(G50.4), and the local algebra $B_\ell$ with its
weighted presentation.  The normalizer conditions were
proved in the first step, all local-condition fields in
the second, and obstruction finiteness in (G50.9).
There is just one framed coefficient-prime local place;
its distinguished comparison matrix is fixed to $1$.
The theorem constructs an actual universal framed object
and a relative weighted presentation with the number of
extra variables of weight $i$ equal to $\dim H^1_c(M_i)$
and the number of equations equal to
$\dim\operatorname{Ob}(M_i)$.
Its coordinates are embedded matrix-entry differences
from the specified basepoint; the extra comparison-matrix
coordinates vanish because that one matrix is $1$.

By the first and second steps, this universal framed
object and $B_\chi$ represent exactly the same functor:
the same $GG_3$ lifts, same killing conditions and same
actual local map killing $P_\chi$.
Their universal properties supply maps in both directions;
each composite specializes the universal representation
to itself and so is the identity by uniqueness.
Transport the constructed presentation along this
isomorphism.  Its relative relation ideal, entry
homogeneity and coordinates are preserved, and (G50.9)
gives exactly (G50.2).  The coordinate weights are
matrix-weight differences, hence in $[-2,2]$.
This is the complete universal-object comparison used
by \path{UnitaryThreeGroups/BoxAbs.lean},
\texttt{boxPres\_abs}, with the verified arbitrary cubic
basepoint in place of its specialized wrapper.
\end{proof}

\begin{proposition}[AC-G50-normalization: the unchanged eight equations]
\label{AC-G50-normalization}
The normalized ring of AC-G49 is the quotient of $B_\chi$
by the eight actual homogeneous equations
\[
 \begin{split}
 &r(c)_{00}-1,\quad r(c)_{11}-1,\quad r(c)_{01},\quad r(c)_{10},\\
 &r(c)_{02},\quad r(c)_{20},\quad z(c,g_1)_{01},\quad z(c,g_1)_{10},
 \end{split}                                                  \tag{G50.10}
\]
Here every $r(c)_{ab}$ denotes the corresponding matrix-component
entry.
with respective weights
$(0,0,1,-1,2,-2,1,-1)$.
Accordingly, after lifting these equations to the ambient
relative presentation, the total equation counts by
weights $(-2,-1,0,1,2)$ are
\[
                 (a_{-2}+1,\ a_{-1}+3,\ a_0+3,
                                      \ a_1+3,\ a_2+1).      \tag{G50.11}
\]
The basepoint kills every equation.  The unit used in the
source marked-element normalization remains a unit:
\[
                z(c,g_1)_{00}\big|_{r_\chi}
                   =q(g_1)^j\bigl(\gamma(g_1)-\gamma(g_1)^{-1}\bigr).
                                                                  \tag{G50.12}
\]
\end{proposition}
\begin{proof}
Since $c^2=1$ and the multiplier at $c$ is $-1$, the
matrix $A=r(c).\mathrm{mat}$ satisfies
$A=JA^tJ$.  Thus $A_{12}=A_{01}$,
$A_{21}=A_{10}$ and $A_{22}=A_{00}$.
The first six equations in (G50.10) are therefore exactly
$A=I$.  The last two are exactly the remaining
\texttt{Normalised} fields.
Universality, or the quotient proof of AC-G49, identifies
the resulting quotient with $R_\chi$.
This proof is independent of the original residual seed;
it is the full \texttt{norm\_iff} argument in
\path{UnitaryThreeGroups/ExistsSurjectionUnivNormCondKerSpanNEq.lean}.

Each entry has its matrix weight, constants have weight
zero, and the two \texttt{zmat} differences have their
entry weights.  This proves the displayed list.  Add its
multiplicities to (G50.2) to obtain (G50.11).
At the basepoint the mixed entries vanish and $r_\chi(c)=I$.
Finally $\eta(g_1)=1$, and anticyclotomy gives
$q(cg_1^{-1}c)=q(g_1)$, whereas
$\gamma(cg_1^{-1}c)=\gamma(g_1)^{-1}$.
This proves (G50.12).  Its first factor and
$\gamma(g_1)$ are units, and the original marked
condition says $\gamma(g_1)^2-1$ is a unit.
Thus the modified marked gap is still a unit, with no
condition $\chi(g_1)=1$.
\end{proof}

\paragraph{What this closes and what it does not.}
At the coefficient primes covered by AC-G49-inert,
AC-G50 supplies the actual relative relation count
previously missing from the mere weighted presentation.
The generic construction uses the source's full
\texttt{LemmaP}, not an assertion that its special
\texttt{Dat} record already accepts the cubic character.
The next conversion to the complete \texttt{Model}
record must still check its entry-density and integral
base-derivative fields against the changed basepoint.
No Type I automorphy, zero mixed tangent image, or numerical
Selmer upper bound is asserted here.

\paragraph{Geometric dictionary.}
The finite residual representations are finite lisse local
systems under the previously audited SGA~1, Expos\'e~V,
\S\S4 and~7 dictionary.  Their framed lifts form the
explicit deformation functor used here; the local conditions
cut out closed subfunctors.  The groups $H^1_c(M_i)$ and
$\operatorname{Ob}(M_i)$ count the variable and relation
spaces in the retained \texttt{LemmaP}.  The existence of
that presentation and the arithmetic Euler-characteristic
identity are attributed to the cited deformation and
cohomology inputs, not to SGA~1.  No additional archival
locator or geometric realization is claimed.
